% =============================================================================
% Project : THE_PHILOSOPHY_OF_VIABILITY
% Book    : Book I — Ontology of Viability
% Book ID : VS.B01
% Part    : Part I — Identity and Manifestation
% Part ID : VS.B01.P01
% Chapter : Chapter 08 — Comparative Philosophical Examination
% Chapter ID : VS.B01.P01.C08
%
% File    :
%   2_MAINMATTER/PART_I_IDENTITY_AND_MANIFESTATION/
%   CHAPTER_08_COMPARATIVE_PHILOSOPHICAL_EXAMINATION/
%   sections/section_17_mulla_sadra.tex
%
% Section : 17 — Mulla Sadra
% Section ID : VS.B01.P01.C08.S17
%
% Purpose :
%   Reconstruct Mulla Sadra's metaphysics of existence, quiddity,
%   gradation, individuation, and substantial motion, and examine
%   its relation to the foundational distinctions of Being and
%   Identity in Book I.
%
% =============================================================================

\section{Mulla Sadra}
\label{sec:ch08_s17}


% =============================================================================
% SS01 — SCOPE AND COMPARATIVE QUESTION
% =============================================================================

\subsection{Scope and Comparative Question}
\label{subsec:ch08_s17_ss01}

% Structural ID: VS.B01.P01.C08.S17.SS01.P001
\phantomsection
\label{para:VS.B01.P01.C08.S17.SS01.P001}

The examination of Mulla Sadra introduces a metaphysical
position in which existence is treated as the fundamental
reality of what exists, while quiddity, individuation,
gradation, and substantial motion are understood through
their relation to existence. The comparison is especially
significant because Sadra does not merely investigate
the persistence of an entity through accidental change.
He questions the assumption that the existent is a
substance that remains ontologically static while
successive alterations occur around it. His account
places becoming within the very actuality of existence.

The present comparison concerns the relation between
this metaphysics and the ontological distinctions
established in Part I. Its central question is whether
the primacy and modulation of existence, together with
the principle of individuation and substantial motion,
provide an account of the existent as the existent
that it is, or whether a further question concerning
Identity remains philosophically necessary.

% Structural ID: VS.B01.P01.C08.S17.SS01.P002
\phantomsection
\label{para:VS.B01.P01.C08.S17.SS01.P002}

The comparison therefore addresses five connected
problems. First, what does Sadra mean by the reality
and primacy of existence? Second, how does the
gradation of existence explain the unity and
multiplicity of existents? Third, in what sense does
existence function as the principle of individuation?
Fourth, how does substantial motion alter the
metaphysical understanding of the subject that
undergoes change? Fifth, does the Sadrian account
render the distinction between Being and Identity
unnecessary, or does it leave an ontological question
that cannot be answered by existence and individuation
alone?

% Structural ID: VS.B01.P01.C08.S17.SS01.P003
\phantomsection
\label{para:VS.B01.P01.C08.S17.SS01.P003}

The comparison must begin by distinguishing the
historical doctrines of Sadra from the interpretation
offered by later scholarship and from the conclusions
developed here. The terms \textit{aṣālat al-wujūd},
\textit{tashkīk al-wujūd}, \textit{tashakhkhuṣ}, and
\textit{al-ḥaraka al-jawhariyya} belong to Sadra's
philosophical vocabulary. Their relation to Being,
Identity, Determination, Invariance, and Continuity
in Book I must be argued rather than presumed.

% Structural ID: VS.B01.P01.C08.S17.SS01.P004
\phantomsection
\label{para:VS.B01.P01.C08.S17.SS01.P004}

The inquiry is consequently comparative rather than
genealogical or justificatory. Sadra's authority does
not establish the truth of Book I, and the conceptual
architecture of Book I does not determine in advance
what Sadra must have meant. The purpose is to identify
the strongest point of contact between the two
positions, formulate the strongest challenge each
raises against the other, and determine whether
their concepts can withstand a test of substitution
without losing their respective meanings.

\index[name]{Mulla Sadra}
\index[concept]{Primacy of existence}
\index[concept]{Substantial motion}

% =============================================================================
% SS02 — HISTORICAL AND TEXTUAL LOCATION
% =============================================================================

\subsection{Historical and Textual Location}
\label{subsec:ch08_s17_ss02}

% Structural ID: VS.B01.P01.C08.S17.SS02.P001
\phantomsection
\label{para:VS.B01.P01.C08.S17.SS02.P001}

\.{S}adr al-D\={i}n Mu\.{h}ammad ibn Ibr\={a}h\={i}m Qaw\={a}m\={i} Sh\={i}r\={a}z\={i},
known as Mulla Sadra and also as \.{S}adr al-Muta'allih\={i}n,
was born in Shiraz around 1571--1572 and died in
1635--1636. His intellectual formation took place
within the philosophical and religious culture of
Safavid Iran. He studied with Mir D\={a}m\={a}d and
Shaykh Bah\={a}' al-D\={i}n al-'\={A}mil\={i}, engaging with
Peripatetic philosophy, Illuminationist thought,
the philosophical legacy of Avicenna, and the
scriptural sciences of the Twelver Shi'i tradition.
His work cannot therefore be adequately represented
as an isolated reaction against Aristotelian
metaphysics or as a simple continuation of one
earlier school. It developed through a sustained
reconstruction of several philosophical and
theological inheritances.\footnote{For the historical
and intellectual setting, see Sajjad H. Rizvi,
``Mulla Sadra,'' \textit{The Stanford Encyclopedia
of Philosophy}, Spring 2021 edition, sections 1
and 2; and Rizvi, ``Moll\={a} \.{S}adr\={a} Sh\={i}r\={a}zi,''
\textit{Encyclopaedia Iranica}.}

% Structural ID: VS.B01.P01.C08.S17.SS02.P002
\phantomsection
\label{para:VS.B01.P01.C08.S17.SS02.P002}

The principal work of Sadra's mature metaphysics
is \textit{al-\.{H}ikma al-muta'\={a}liya f\={i} al-asf\={a}r
al-'aqliyya al-arba'a}, commonly known as
\textit{al-Asf\={a}r al-arba'a} or \textit{The Four
Journeys}. It is a large philosophical and
theological summa whose organization brings
together metaphysics, theology, natural philosophy,
psychology, and eschatology. Other works relevant
to the present comparison include
\textit{al-Shaw\={a}hid al-rub\={u}biyya},
\textit{al-Mabda' wa-al-Ma'\={a}d},
\textit{al-\.{H}ikma al-'Arshiyya}, and
\textit{Kit\={a}b al-Mash\={a}'ir}. These works should
not be treated as interchangeable summaries.
Their different purposes and argumentative
contexts must be considered when a particular
doctrine is reconstructed.\footnote{For the
bibliography and editorial history of Sadra's
works, see Rizvi, ``Mulla Sadra,'' sections
1.2 and bibliography. The nine-volume critical
edition of \textit{al-Asf\={a}r} was published
between 2001 and 2005 under the general
editorship of Sayyid Muhammad Khaminihi.}

% Structural ID: VS.B01.P01.C08.S17.SS02.P003
\phantomsection
\label{para:VS.B01.P01.C08.S17.SS02.P003}

The title \textit{al-ḥikma al-muta'āliya} is
often translated as ``transcendent philosophy''
or ``transcendent wisdom.'' Neither translation
should be mistaken for a complete definition of
Sadra's philosophical method. His inquiry
combines demonstrative reasoning with forms
of intellectual and spiritual realization, and
it places philosophical argument in relation
to scriptural interpretation. The relation
between demonstration and intuition is not
adequately described as the replacement of
reason by mystical experience. Sadra argues
for their compatibility within a broader
conception of philosophical wisdom.\footnote{See
Rizvi, ``Mulla Sadra,'' section 2.2; and
Ibrahim Kalin, \textit{Knowledge in Later
Islamic Philosophy: Mulla Sadra on Existence,
Intellect, and Intuition} (Oxford: Oxford
University Press, 2010), especially the
discussion of existence and intelligibility.}

% Structural ID: VS.B01.P01.C08.S17.SS02.P004
\phantomsection
\label{para:VS.B01.P01.C08.S17.SS02.P004}

The earlier comparative material in the
project sometimes presents Sadra through a
short account of three worlds: intellect,
image, and nature. This account may identify
one relevant cosmological arrangement, but
it cannot serve as a complete reconstruction
of his metaphysics. The central doctrines
examined here concern the reality of existence,
its modulation, the individuation of existents,
and substantial motion. Their significance
cannot be reduced to a fixed enumeration of
cosmological regions or transferred directly
to the proposed realms and hubs of Book I.

% Structural ID: VS.B01.P01.C08.S17.SS02.P005
\phantomsection
\label{para:VS.B01.P01.C08.S17.SS02.P005}

The historical reconstruction must also
distinguish Sadra's own formulations from
later Sadrian systematization. The philosophical
school that developed after him elaborated,
organized, and sometimes reformulated aspects
of his thought. A doctrine found in a later
commentator may illuminate Sadra's system,
but it cannot automatically be attributed
to Sadra in precisely the same form. The
present examination therefore uses later
scholarship to clarify the interpretive
problems while retaining the distinction
between Sadra's texts and their subsequent
reception.

% Structural ID: VS.B01.P01.C08.S17.SS02.P006
\phantomsection
\label{para:VS.B01.P01.C08.S17.SS02.P006}

The historical and textual position established
here is sufficient for the comparison that
follows. It does not require a complete history
of the School of Isfahan, the later Sadrian
tradition, or the entire reception of
\textit{al-ḥikma al-muta'āliya}. Those matters
remain relevant where they clarify a disputed
interpretation, but they do not replace the
reconstruction of the arguments concerning
existence, individuation, and change.

\index[concept]{Transcendent wisdom}


% =============================================================================
% SS03 — THE ARCHITECTURE OF TRANSCENDENT WISDOM
% =============================================================================

\subsection{The Architecture of Transcendent Wisdom}
\label{subsec:ch08_s17_ss03}

% Structural ID: VS.B01.P01.C08.S17.SS03.P001
\phantomsection
\label{para:VS.B01.P01.C08.S17.SS03.P001}

Sadra's metaphysics is organized around the
question of existence. This does not mean
that every philosophical problem is reduced
to a single formula or that all distinctions
within his system become dispensable. Rather,
the reality of existence provides the
fundamental orientation through which
quiddity, causation, unity, multiplicity,
individuation, motion, and the relation
between the divine principle and contingent
existents are reconsidered.

% Structural ID: VS.B01.P01.C08.S17.SS03.P002
\phantomsection
\label{para:VS.B01.P01.C08.S17.SS03.P002}

The architecture of Sadra's philosophy
cannot be understood as a mere combination
of three independent traditions. Peripatetic
demonstration, Illuminationist insight, and
Sufi metaphysical resources are important
to his intellectual formation, but his
philosophy does not simply place their
doctrines alongside one another. It
reconstructs their questions within a
distinctive account of existence. The
result is a system in which logical
demonstration, intellectual presence,
and scriptural interpretation have
interrelated but non-identical roles.

% Structural ID: VS.B01.P01.C08.S17.SS03.P003
\phantomsection
\label{para:VS.B01.P01.C08.S17.SS03.P003}

The first journey of \textit{al-Asfār}
establishes philosophical and metaphysical
principles; the subsequent journeys examine
the divine principle, the relation between
God and the world, and the human soul and
its ultimate condition. This arrangement
is important because it situates the
doctrine of existence within a wider
metaphysical inquiry. The primacy of
existence is not an isolated thesis about
the meaning of a word. It concerns the
ontological constitution of contingent
realities and their dependence upon the
Necessary.

% Structural ID: VS.B01.P01.C08.S17.SS03.P004
\phantomsection
\label{para:VS.B01.P01.C08.S17.SS03.P004}

The present section concentrates on the
ontological consequences of this architecture.
The doctrines of knowledge by presence
and the union of intellect and intelligible
will be examined only insofar as they
clarify the relation between existence
and knowing. Their epistemological
implications are not imported into the
foundations of Book I. The comparative
question remains ontological: what is
the reality of the existent, what makes
it determinate, and how is its actuality
related to change?

% Structural ID: VS.B01.P01.C08.S17.SS03.P005
\phantomsection
\label{para:VS.B01.P01.C08.S17.SS03.P005}

The architecture of Sadra's system thus
provides a demanding interlocutor for
Book I. It does not merely offer an
alternative vocabulary for identity.
It proposes that the reality which
grounds the existent, differentiates
existents, and undergoes substantial
motion is existence itself. The
comparison must determine whether
this proposal answers the same
ontological questions as the distinction
between Being and Identity, or whether
the two inquiries operate at different
levels.

\index[concept]{Existence}


% =============================================================================
% SS04 — EXISTENCE AND QUIDDITY
% =============================================================================

\subsection{Existence and Quiddity}
\label{subsec:ch08_s17_ss04}

% Structural ID: VS.B01.P01.C08.S17.SS04.P001
\phantomsection
\label{para:VS.B01.P01.C08.S17.SS04.P001}

The distinction between existence
(\textit{wujūd}) and quiddity
(\textit{māhiyya}) is indispensable
to Sadra's metaphysics. Quiddity concerns
what a thing is, or the intelligible
content through which it can be defined
and distinguished conceptually. Existence
concerns that the thing is and, more
fundamentally, the concrete actuality
through which it is present as an existent.
The distinction does not establish two
independently existing components
within every contingent entity. It
articulates two different aspects of
the analysis of what exists.

% Structural ID: VS.B01.P01.C08.S17.SS04.P002
\phantomsection
\label{para:VS.B01.P01.C08.S17.SS04.P002}

Sadra rejects the assumption that
quiddity is the independently real
foundation to which existence is
subsequently attached. The quiddity
of a contingent thing does not first
subsist as an actual entity awaiting
the addition of existence. Rather,
the existent is actual through existence,
and its quiddity is understood in
relation to that actuality. The
ontological priority assigned to
existence is therefore directed
against an account in which the
essence of a contingent thing is
treated as the primary reality
and existence as an additional
feature.\footnote{For Sadra's discussion
of the relation between existence
and essence in contingent beings,
see \textit{al-Asfār}, vol. I,
pp. 289--292, as discussed by Rizvi,
``Mulla Sadra,'' section 3.1.}

% Structural ID: VS.B01.P01.C08.S17.SS04.P003
\phantomsection
\label{para:VS.B01.P01.C08.S17.SS04.P003}

This position requires a distinction
between quiddity as a conceptual
determination and existence as
concrete ontological reality. The
concept of existence may be grasped
intellectually, but the reality of
existence is not exhausted by the
concept through which it is
understood. Sadra's discussion
therefore distinguishes existence
as a concept from existence as
an extra-mental reality. Confusing
these two senses would reduce
the ontological question to a
logical or semantic analysis.

% Structural ID: VS.B01.P01.C08.S17.SS04.P004
\phantomsection
\label{para:VS.B01.P01.C08.S17.SS04.P004}

The distinction between existence
and quiddity must not be translated
directly into the distinction between
Being and Identity in Book I.
Quiddity is not equivalent to Identity,
and the fact that Sadra denies
independent ontological primacy
to quiddity does not establish
that he denies every question
concerning the identity of an existent.
Book I's Identity concerns the
existent as the existent that it is.
It is not a conceptual inventory
of defining characteristics and
cannot be treated as a synonym
for a thing's quiddity.

% Structural ID: VS.B01.P01.C08.S17.SS04.P005
\phantomsection
\label{para:VS.B01.P01.C08.S17.SS04.P005}

The comparison must therefore
preserve two distinct relations.
Sadra examines the ontological
priority of existence over quiddity.
Book I examines the analytical
distinction and ontological
inseparability of Being and Identity,
together with the absence of
ontological priority between them.
These are not the same pair of
concepts. Any conclusion that
Sadra's primacy of existence
directly contradicts the foundational
relation between Being and Identity
would depend upon an unargued
identification of Identity with
quiddity.

\index[concept]{Quiddity}
\index[concept]{Existence and quiddity}


% =============================================================================
% SS05 — THE PRIMACY OF EXISTENCE
% =============================================================================

\subsection{The Primacy of Existence}
\label{subsec:ch08_s17_ss05}

% Structural ID: VS.B01.P01.C08.S17.SS05.P001
\phantomsection
\label{para:VS.B01.P01.C08.S17.SS05.P001}

The doctrine of the primacy of existence
(\textit{aṣālat al-wujūd}) maintains that
existence, rather than quiddity, is
the fundamental ontological reality.
The doctrine is not merely the assertion
that things possess existence. It
concerns the status of existence
as the actuality through which
contingent realities are real.
Quiddities do not constitute an
independent ontological domain
that subsequently receives existence.
Their actuality is understood
through the existence in which
they are realized.

% Structural ID: VS.B01.P01.C08.S17.SS05.P002
\phantomsection
\label{para:VS.B01.P01.C08.S17.SS05.P002}

The doctrine also transforms the
interpretation of contingency.
A contingent existent is not
an independently subsisting essence
to which existence is externally
added. Its dependence concerns
the very actuality through which
it exists. This enables Sadra to
place the relation between the
Necessary and contingent existents
within an existential rather than
an essence-first account. The
causal dependence of the contingent
is therefore not an accidental
relation between two already
complete realities; it concerns
the dependence of its existence
upon its principle.

% Structural ID: VS.B01.P01.C08.S17.SS05.P003
\phantomsection
\label{para:VS.B01.P01.C08.S17.SS05.P003}

The primacy of existence must
not be confused with the claim
that the concept of existence
is a universal genus shared
by individual things in the
manner of a common species.
Sadra's analysis distinguishes
the simplicity of the concept
from the concrete reality to
which it refers. Existence is
not an ordinary genus whose
instances are differentiated
by external specific differences.
Its predication and its
ontological realization require
a more complex account, developed
through the doctrine of
\textit{tashkīk}.

% Structural ID: VS.B01.P01.C08.S17.SS05.P004
\phantomsection
\label{para:VS.B01.P01.C08.S17.SS05.P004}

The primacy of existence consequently
establishes the starting point
for Sadra's account of unity,
multiplicity, and individuation.
If existence is the fundamental
reality, then the differentiation
of existents must be explained
through the modes, degrees,
or intensities of existence
rather than through a plurality
of independently real quiddities.
The next question is therefore
not simply why many things have
the same concept of existence,
but how one existential reality
can be differentiated without
being reduced to a collection
of unrelated substances.

% Structural ID: VS.B01.P01.C08.S17.SS05.P005
\phantomsection
\label{para:VS.B01.P01.C08.S17.SS05.P005}

The primacy of existence is
a foundational commitment within
Sadra's own metaphysical system.
Its status in the present comparison
is different. It is a historically
reconstructed thesis whose
interpretation is supported
by primary-text scholarship.
It is not a foundational commitment
adopted by Book I merely because
it offers a powerful account
of the actuality of existents.

\index[concept]{Primacy of existence}


% =============================================================================
% SS06 — MODULATION AND GRADATION
% =============================================================================

\subsection{The Modulation and Gradation of Existence}
\label{subsec:ch08_s17_ss06}

% Structural ID: VS.B01.P01.C08.S17.SS06.P001
\phantomsection
\label{para:VS.B01.P01.C08.S17.SS06.P001}

The doctrine of modulation and gradation
(\textit{tashkīk al-wujūd}) develops
the primacy of existence by explaining
how existence can be one reality
while being realized in differentiated
degrees. The differences between
existents are not treated as
differences between independent
instances of an ordinary universal.
They concern the existential
reality itself, which admits
of degrees of intensity,
perfection, priority, and dependence.

% Structural ID: VS.B01.P01.C08.S17.SS06.P002
\phantomsection
\label{para:VS.B01.P01.C08.S17.SS06.P002}

The doctrine thus addresses the
relation between unity and
multiplicity. Existence is not
multiplied into unrelated realities
merely because many existents
are encountered. Nor is the
multiplicity of existents dismissed
as an illusion without ontological
significance. The degrees are
real, and their differences
are articulated through the
intensification or diminution
of existence. Unity and multiplicity
are therefore understood within
one modulated existential reality.

% Structural ID: VS.B01.P01.C08.S17.SS06.P003
\phantomsection
\label{para:VS.B01.P01.C08.S17.SS06.P003}

The gradation of existence is
not equivalent to a hierarchy
of conceptual classifications.
Its differentiations are intended
to be ontological. The relation
between more intense and less
intense existence is not simply
a relation between more and
less extensive descriptions,
nor does it depend upon the
number of predicates attributed
to a thing. It concerns the
mode of actuality through which
an existent is what it is.
The precise interpretation of
this modulation, however, has
been the subject of philosophical
discussion and should not be
reduced to a single schematic
formula.\footnote{For the account
of modulation and gradation,
see Rizvi, ``Mulla Sadra,''
section 3.2, and Rizvi,
\textit{Mulla Sadra and Metaphysics:
The Modulation of Being}.}

% Structural ID: VS.B01.P01.C08.S17.SS06.P004
\phantomsection
\label{para:VS.B01.P01.C08.SS06.P004}

The distinction between gradation
and Invariance is essential
for the present comparison.
Invariance, as examined in
Book I, requires specification
of a transformation, a domain,
a respect, and the limits
within which something remains
invariant. Sadra's gradation
concerns differentiated degrees
of existence. It does not,
without an additional argument,
identify what remains invariant
under a specified transformation.
A graded existential reality
may undergo substantial change;
the existence of degrees does
not itself establish an invariant
identity across them.

% Structural ID: VS.B01.P01.C08.S17.SS06.P005
\phantomsection
\label{para:VS.B01.P01.C08.S17.SS06.P005}

Nor should the gradation of
existence be identified with
the proposed realms and hubs
of Book I. Sadra's hierarchy
is an account of degrees of
existence. The proposed realms
and hubs concern a different
architecture and have not been
established as degrees of
existential intensity. A
structural resemblance between
stratified accounts of reality
cannot establish that their
constituent domains have the
same ontological function.

\index[concept]{Modulation of existence}
\index[concept]{Gradation of existence}


% =============================================================================
% SS07 — UNITY AND MULTIPLICITY
% =============================================================================

\subsection{Unity and Multiplicity in Existence}
\label{subsec:ch08_s17_ss07}

% Structural ID: VS.B01.P01.C08.S17.SS07.P001
\phantomsection
\label{para:VS.B01.P01.C08.S17.SS07.P001}

Sadra's account of unity and multiplicity
is inseparable from the primacy
and gradation of existence. The
one and the many are not treated
as two mutually exclusive descriptions
between which metaphysics must
choose. The unity of existence
and the multiplicity of existents
are understood through the
differentiated realization of
existence itself. The degrees
are not independent substances
that happen to share a common
name; neither are they merely
subjective divisions imposed
upon an otherwise undifferentiated
concept.

% Structural ID: VS.B01.P01.C08.S17.SS07.P002
\phantomsection
\label{para:VS.B01.P01.C08.S17.SS07.P002}

The expression \textit{basīṭ
al-ḥaqīqa kull al-ashyā'}—
the simple reality is all things—
belongs to this broader metaphysical
orientation. It concerns the
relation between the simplicity
of the ultimate principle and
the multiplicity of what depends
upon it. It should not be treated
as a simple assertion that every
finite existent is numerically
identical with God, or that
all distinctions among existents
are meaningless. The interpretation
of simplicity, dependence,
and multiplicity requires
attention to the different
levels and modes of Sadra's
argument.

% Structural ID: VS.B01.P01.C08.S17.SS07.P003
\phantomsection
\label{para:VS.B01.P01.C08.S17.SS07.P003}

The Sadrian account is therefore
a serious alternative to
metaphysical pluralism understood
as the independent subsistence
of many unrelated realities.
It seeks to explain how
differentiation can be real
without requiring that existence
be divided into separate
existential principles. The
doctrine of gradation performs
the principal explanatory work
in this respect: multiplicity
is articulated through degrees
of one existential reality.

% Structural ID: VS.B01.P01.C08.S17.SS07.P004
\phantomsection
\label{para:VS.B01.P01.C08.S17.SS07.P004}

The comparison with Book I
must nevertheless preserve
the distinction between unity
and Identity. Book I does
not define Identity as the
unity of an existent, nor
does it assume that the unity
of one existential reality
settles the question of how
a particular existent is
the existent that it is.
Sadra's account of unity
and multiplicity may explain
the relation between degrees
of existence, but the further
question concerning the
ontological specificity of
each existent remains relevant
to the present inquiry.

% Structural ID: VS.B01.P01.C08.S17.SS07.P005
\phantomsection
\label{para:VS.B01.P01.C08.S17.SS07.P005}

The doctrine also places pressure
upon any account that treats
multiplicity as sufficient
evidence for a plurality
of independent identities.
The existence of many
differentiated existents does
not by itself establish that
each difference corresponds
to an independently constituted
Identity in the sense of
Book I. Conversely, the unity
of existence does not by itself
establish that there is only
one Identity. The relation
between existential multiplicity
and the possible plurality
of identities remains a
separate ontological problem.

\index[concept]{Unity and multiplicity}
\index[concept]{Simple reality}


% =============================================================================
% SS08 — EXISTENCE AND INDIVIDUATION
% =============================================================================

\subsection{Existence and the Principle of Individuation}
\label{subsec:ch08_s17_ss08}

% Structural ID: VS.B01.P01.C08.S17.SS08.P001
\phantomsection
\label{para:VS.B01.P01.C08.S17.SS08.P001}

One of the most important consequences
of Sadra's metaphysics concerns
individuation. Existence is not
only the source of the actuality
of an existent; it also functions
as the principle through which
the existent is individuated.
The term \textit{tashakhkhuṣ}
refers to the constitution
or determination of an existent
as an individual. Sadra's position
therefore gives existence a
role in explaining why an
existent is this individual
rather than merely an instance
of a definable quiddity.

% Structural ID: VS.B01.P01.C08.S17.SS08.P002
\phantomsection
\label{para:VS.B01.P01.C08.S17.SS08.P002}

This doctrine is especially
important because it places
individuation within the reality
of existence rather than
assigning it exclusively to
matter, a bundle of properties,
or an external differentiating
condition. The individual is
not first constituted as
a complete quiddity and then
made numerically distinct
through an additional feature.
Its individuation is grounded
in its existential actuality.
Rizvi reports Sadra's formulation
that existence is the principle
of individuality and the source
of the quiddity of the thing.\footnote{See
\textit{al-Asfār}, vol. III,
p. 351, as cited and discussed
in Rizvi, ``Mulla Sadra,''
section 4.5. The precise
interpretation of the passage
must be distinguished from
the broader claim that every
use of ``identity'' in later
commentary is equivalent
to Book I's Identity.}

% Structural ID: VS.B01.P01.C08.S17.SS08.P003
\phantomsection
\label{para:VS.B01.P01.C08.S17.SS08.P003}

The principle of individuation
is not, however, automatically
equivalent to Identity. Individuation
concerns the constitution
of an existent as an individual.
Identity in Book I concerns
the existent as the existent
that it is. The two questions
are closely related because
an individual existent is
the bearer of an identity
question. Yet the constitution
of an individual and the
ontological meaning of its
Identity are not interchangeable
without further argument.

% Structural ID: VS.B01.P01.C08.S17.SS08.P004
\phantomsection
\label{para:VS.B01.P01.C08.S17.SS08.P004}

The distinction becomes especially
important where the possibility
of multiple identities is
considered. A theory of individuation
may explain how an existent
is distinguished as one individual.
It does not thereby determine
whether that individual can
involve more than one identity,
whether its identity is exhausted
by its individuality, or whether
the relations and manifestations
of one individual can involve
ontologically distinct identities.
These questions remain open
within Book I and cannot
be resolved by importing
the Sadrian principle of
individuation as though it
were already a complete theory
of Identity.

% Structural ID: VS.B01.P01.C08.S17.SS08.P005
\phantomsection
\label{para:VS.B01.P01.C08.S17.SS08.P005}

Sadra's account nevertheless
poses a decisive challenge.
If existence itself performs
the work of individuation,
then Book I must explain
what additional or distinct
ontological work is performed
by Identity. The answer cannot
be that Identity is simply
another name for individuation,
nor can it be that Identity
is an independently existing
entity added to the existent.
The issue is whether the
analytical distinction between
Being and Identity clarifies
an aspect of ontological
actuality that the principle
of individuation does not
exhaust.

\index[concept]{Individuation}
\index[concept]{Tashakhkhus}


% =============================================================================
% SS09 — SUBSTANTIAL MOTION
% =============================================================================

\subsection{Substantial Motion}
\label{subsec:ch08_s17_ss09}

% Structural ID: VS.B01.P01.C08.S17.SS09.P001
\phantomsection
\label{para:VS.B01.P01.C08.S17.SS09.P001}

The doctrine of substantial motion
(\textit{al-ḥaraka al-jawhariyya})
extends the dynamism of existence
to the level of substance itself.
In the Aristotelian account,
motion is commonly analyzed
through changes in quantity,
quality, place, or other
accidental determinations while
the substance that undergoes
those changes remains the
underlying subject. Sadra
challenges the assumption that
substance is ontologically
static in this manner. The
change of an existent is
not confined to alterations
in accidents; the substance
itself undergoes motion.

% Structural ID: VS.B01.P01.C08.S17.SS09.P002
\phantomsection
\label{para:VS.B01.P01.C08.S17.SS09.P002}

Substantial motion is therefore
not simply the claim that
substances move. It is the
claim that the existential
actuality of the substance
is itself dynamic. The existent
is not a fixed substrate
to which a succession of
external changes is attached.
Its actuality unfolds through
a continuous process of
substantial becoming. This
doctrine is connected to
the primacy and gradation
of existence because the
changing existent is understood
through the modulated reality
of existence rather than
through a permanently fixed
quiddity.\footnote{For the
distinction between substantial
motion and the ordinary motion
of substances, see Fazlur Rahman,
\textit{The Philosophy of Mulla
Sadra}, pp. 95--108; and Rizvi,
``Mulla Sadra,'' section 3.2.}

% Structural ID: VS.B01.P01.C08.S17.SS09.P003
\phantomsection
\label{para:VS.B01.P01.C08.S17.SS09.P003}

The doctrine changes the formulation
of the problem of change.
The question is no longer
only how a substance can
remain numerically the same
while its accidents alter.
It also concerns how an
existent can be understood
as one developing actuality
when its very substantial
existence is in motion.
The subject of change cannot
be explained simply as an
unchanging substrate standing
outside the process. The
account must explain the
relation between the actuality
that changes and the continuity
through which the change
is attributed to an existent.

% Structural ID: VS.B01.P01.C08.S17.SS09.P004
\phantomsection
\label{para:VS.B01.P01.C08.SS09.P004}

Substantial motion must not
be identified with Expression
in Book I. Expression concerns
a distinct ontological problem
within the account of Identity
and its manifestation. Substantial
motion concerns the transformation
of the existent at the level
of its substance. The two
doctrines may be compared
because both challenge static
accounts of reality, but
they do not perform the
same explanatory function.
Nor can substantial motion
be reduced to a change in
Manifestation without first
demonstrating that the
substantial transformation
described by Sadra is
ontologically equivalent
to the relation examined
in Book I.

% Structural ID: VS.B01.P01.C08.S17.SS09.P005
\phantomsection
\label{para:VS.B01.P01.C08.S17.SS09.P005}

The significance of substantial
motion for Book I is therefore
not that it supplies a ready-made
solution to the problem of
Identity through change.
It intensifies that problem.
If the existent changes
in its substantial actuality,
the account must explain
whether the same individual
continues, whether successive
existential stages constitute
one developing individual,
and what grounds the attribution
of continuity across those
stages. These questions cannot
be settled by merely naming
the process substantial motion.

\index[concept]{Substantial motion}
\index[concept]{Haraka jawhariyya}


% =============================================================================
% SS10 — THE SUBJECT OF MOTION AND CONTINUITY
% =============================================================================

\subsection{The Subject of Motion and the Problem of Continuity}
\label{subsec:ch08_s17_ss10}

% Structural ID: VS.B01.P01.C08.S17.SS10.P001
\phantomsection
\label{para:VS.B01.P01.C08.S17.SS10.P001}

Substantial motion raises a question
that is central to the comparison:
what is the subject of change
when the substance itself is
in motion? If the existent
is not a fixed substrate
that remains wholly outside
the process, then continuity
cannot be explained merely
by referring to an unaltered
substance beneath successive
states. The subject must
be understood through the
existential process in which
its actuality is realized.

% Structural ID: VS.B01.P01.C08.S17.SS10.P002
\phantomsection
\label{para:VS.B01.P01.C08.S17.SS10.P002}

Rizvi's reconstruction draws
attention to the tension between
the continuous transformation
of the existent and the
recognition of an individual
as the same individual.
In one discussion, Sadra's
account is presented through
the claim that the individual
remains who he is while
the intensity and mode of
existence change. In another,
the substantial transformation
of the existent is described
in terms that distinguish
the young Zayd from the
old Zayd as successive
existential stages. These
formulations should not be
collapsed into a single
unqualified assertion that
numerical identity is either
always preserved or always
destroyed.\footnote{The discussion
of the individual remaining
``who he is'' refers to
\textit{al-Asfār}, vol. III,
p. 351, as cited by Rizvi.
The discussion of substantial
motion and the changing
existent appears in Rizvi,
``Mulla Sadra,'' section 3.2.
The apparent tension requires
interpretation in relation
to the different questions
being addressed.}

% Structural ID: VS.B01.P01.C08.S17.SS10.P003
\phantomsection
\label{para:VS.B01.P01.C08.S17.SS10.P003}

At least three questions must
therefore be distinguished.
The first concerns the continuity
of an existential process.
The second concerns the
individuality of the subject
through the process. The
third concerns numerical
identity across successive
stages of substantial becoming.
A theory may provide an
account of the continuity
of a process without thereby
establishing that every
successive stage is numerically
identical in every relevant
respect. Likewise, the
recognition of one individual
through change does not
by itself explain the precise
ontological conditions under
which its identity obtains.

% Structural ID: VS.B01.P01.C08.S17.SS10.P004
\phantomsection
\label{para:VS.B01.P01.C08.S17.SS10.P004}

This distinction is important
for Book I because Identity
must not be equated with
persistence, continuity, or
stability. The fact that
an existent continues through
a process does not by itself
define its Identity. Conversely,
a transformation in its
substantial constitution does
not automatically establish
that its Identity has ceased.
The ontological relation
between substantial change
and Identity must be examined
through the constitutive
conditions of the existent,
rather than inferred from
the mere presence or absence
of continuity.

% Structural ID: VS.B01.P01.C08.S17.SS10.P005
\phantomsection
\label{para:VS.B01.P01.C08.S17.SS10.P005}

Sadra's position therefore
offers a powerful challenge
to theories that explain
continuity by presupposing
an entirely static subject.
Yet the challenge does not
remove the need to distinguish
the continuity of becoming,
the individuality of the
existent, and its Identity.
Whether Sadra's account
fully explains their relation
is a question that requires
a closer examination of
his treatment of the soul,
the individual, and the
different stages of existence.

\index[concept]{Continuity through change}


% =============================================================================
% SS11 — TIME AND RENEWAL
% =============================================================================

\subsection{Time, Renewal, and the Actuality of Change}
\label{subsec:ch08_s17_ss11}

% Structural ID: VS.B01.P01.C08.S17.SS11.P001
\phantomsection
\label{para:VS.B01.P01.C08.S17.SS11.P001}

The doctrine of substantial
motion also affects the
understanding of time. Time
is not treated as an entirely
independent container within
which already constituted
substances undergo change.
It is related to the motion
of existence and is understood
as a dimension or measure
of the existential process.
The temporal order is therefore
connected to the actuality
of becoming rather than
being a neutral framework
external to it.

% Structural ID: VS.B01.P01.C08.S17.SS11.P002
\phantomsection
\label{para:VS.B01.P01.C08.S17.SS11.P002}

The renewal of existence
at successive moments follows
from this dynamic account.
The world is not understood
as a collection of substances
whose complete actuality
is established once and
for all. The actuality of
the existent is renewed
through the ongoing process
of substantial motion. This
doctrine is connected to
Sadra's account of creation
and dependence, although
the precise relation between
temporal renewal, divine
causation, and the continuity
of the existent requires
careful reconstruction.

% Structural ID: VS.B01.P01.C08.S17.SS11.P003
\phantomsection
\label{para:VS.B01.P01.C08.S17.SS11.P003}

The renewal of existence
must not be confused with
the conclusion that Identity
is simply replaced at every
moment. That conclusion
would require an additional
argument concerning the
relation between successive
existential stages and the
individual whose development
they constitute. Likewise,
the continuity of the process
does not by itself prove
that every aspect of the
existent remains invariant.
The temporal renewal of
existence and the ontological
question of Identity must
therefore remain analytically
distinct.

% Structural ID: VS.B01.P01.C08.S17.SS11.P004
\phantomsection
\label{para:VS.B01.P01.C08.S17.SS11.P004}

The comparison with Book I
is particularly important
here. Invariance requires
a specified transformation,
domain, respect, type, and
limit. Sadra's account of
temporal renewal concerns
the dynamic actuality of
existence. It does not
automatically identify an
invariant under a defined
class of transformations.
Nor does the absence of
an invariant in one respect
establish the absence of
Identity. The Sadrian account
therefore sharpens the
problem of change without
settling the distinct question
of Invariance as developed
in Book I.

\index[concept]{Time}
\index[concept]{Renewal of existence}


% =============================================================================
% SS12 — THE SOUL AND TRANS-SUBSTANTIAL BECOMING
% =============================================================================

\subsection{The Soul and Trans-Substantial Becoming}
\label{subsec:ch08_s17_ss12}

% Structural ID: VS.B01.P01.C08.S17.SS12.P001
\phantomsection
\label{para:VS.B01.P01.C08.S17.SS12.P001}

Sadra's account of the soul
provides an important case
for examining substantial
motion and individuality.
The soul is not understood
as a fully completed substance
that remains unchanged
throughout its relation to
the body. Its development
is connected to the dynamic
actuality of existence.
The formulation that the
soul is corporeal in its
origination and spiritual
in its subsistence expresses
the relation between its
initial embodiment and
its subsequent existential
development.

% Structural ID: VS.B01.P01.C08.S17.SS12.P002
\phantomsection
\label{para:VS.B01.P01.C08.S17.SS12.P002}

The soul's development is
not merely a succession
of accidental psychological
states. It is treated as
a substantial process in
which the existential actuality
of the soul becomes more
intense and more complete.
The account connects the
formation of the human
soul with its bodily origin
and its possible spiritual
subsistence. The individual
is therefore not explained
solely through the persistence
of a material configuration.
The soul has an ontological
role in the constitution
and continuity of the
individual person.

% Structural ID: VS.B01.P01.C08.S17.SS12.P003
\phantomsection
\label{para:VS.B01.P01.C08.S17.SS12.P003}

This account is relevant
to the problem of individuation
because Sadra does not
treat matter as the sole
principle through which
individual personhood is
constituted. The soul's
existence and development
are central to the explanation
of the individual's identity
and post-mortem condition.
Nevertheless, the particular
Sadrian account of the soul
must not be generalized
into a universal theory
of the individuation of
all existents. Its conclusions
belong to a metaphysics
that includes specific
commitments concerning
the soul, the body, and
eschatological existence.

% Structural ID: VS.B01.P01.C08.S17.SS12.P004
\phantomsection
\label{para:VS.B01.P01.C08.S17.SS12.P004}

The comparison with Book I
must consequently distinguish
the ontological constitution
of a person from the broader
problem of Identity. Sadra's
account may explain how
a human individual develops
through substantial change
and how the soul grounds
its continued existence.
It does not follow that
Identity in Book I should
be defined through the soul,
through personal persistence,
or through a particular
anthropological theory.
The ontological question
of Identity has a broader
scope than the specific
conditions of human personhood.

\index[concept]{Trans-substantial becoming}
\index[concept]{Soul}


% =============================================================================
% SS13 — KNOWLEDGE BY PRESENCE AND THE INTELLECT
% =============================================================================

\subsection{Knowledge by Presence and the Union of Intellect and Intelligible}
\label{subsec:ch08_s17_ss13}

% Structural ID: VS.B01.P01.C08.S17.SS13.P001
\phantomsection
\label{para:VS.B01.P01.C08.S17.SS13.P001}

Sadra's metaphysics is closely
connected to his account
of knowledge. Knowledge by
presence (\textit{al-'ilm
al-ḥuḍūrī}) concerns the
presence of the known to
the knower rather than
the mediation of knowledge
through a mental representation
alone. His account of the
union of the intellect and
the intelligible
(\textit{ittiḥād al-'āqil
wa-al-ma'qūl}) develops
this position by examining
the relation between the
intellect, the act of
intellection, and the
intelligible object.

% Structural ID: VS.B01.P01.C08.S17.SS13.P002
\phantomsection
\label{para:VS.B01.P01.C08.S17.SS13.P002}

These doctrines are not
external additions to Sadra's
metaphysics. They are connected
to his account of existence
and to the claim that
intelligibility is grounded
in existential actuality.
The knower-known relation
is therefore not treated
as a merely external
correspondence between
two independently complete
realities. The relation
between intellect and
intelligible is examined
through the degrees and
modes of existence.

% Structural ID: VS.B01.P01.C08.S17.SS13.P003
\phantomsection
\label{para:VS.B01.P01.C08.S17.SS13.P003}

The doctrine of the union
of intellect and intelligible
is nevertheless a distinct
philosophical problem from
the ontological relation
between Being and Identity.
It concerns the nature
and possibility of knowledge,
including the relation
between the intellecting
subject and the intelligible
content. It must not be
used to establish that
Identity depends upon
cognition, that an existent
requires a knower in order
to be what it is, or that
the ontological actuality
of an entity is constituted
by an epistemic act.

% Structural ID: VS.B01.P01.C08.S17.SS13.P004
\phantomsection
\label{para:VS.B01.P01.C08.S17.SS13.P004}

The relevance of this doctrine
to Book I is therefore
comparative and limited.
It demonstrates how deeply
Sadra integrates the
metaphysics of existence
with the philosophy of
knowledge. It also provides
a useful point of contrast
for the methodological
independence of Book I's
ontological inquiry. The
present comparison does
not deny that cognition
may itself possess an
ontological status. It
maintains that the existence
and Identity of an existent
must not be made dependent
upon the conditions through
which that existent becomes
known.

\index[concept]{Knowledge by presence}
\index[concept]{Union of intellect and intelligible}


% =============================================================================
% SS14 — MULLA SADRA AND IBN ARABI
% =============================================================================

\subsection{Mulla Sadra and Ibn Arabi}
\label{subsec:ch08_s17_ss14}

% Structural ID: VS.B01.P01.C08.S17.SS14.P001
\phantomsection
\label{para:VS.B01.P01.C08.S17.SS14.P001}

Sadra's metaphysics bears
a significant relation to
the philosophical and Sufi
traditions that preceded
him, including the thought
of Ibn Arabi. His account
of existence, unity, and
multiplicity develops within
an intellectual environment
in which Akbarian metaphysical
themes had become important
for philosophical reflection.
Yet the relation between
Sadra and Ibn Arabi cannot
be described as a simple
transmission of one complete
system into another.

% Structural ID: VS.B01.P01.C08.S17.SS14.P002
\phantomsection
\label{para:VS.B01.P01.C08.S17.SS14.P002}

The comparison with Section~\ref{sec:ch08_s16}
is therefore not a repetition
of the earlier examination.
Ibn Arabi's account was
examined through the unity
of the Real, divine names,
self-disclosure, fixed
entities, preparedness,
and the Barzakh. Sadra's
system places the reality
and gradation of existence
at the center of its
philosophical architecture
and develops a distinct
account of individuation
and substantial motion.
The two positions share
a serious concern with
the relation between unity
and multiplicity, but
their conceptual organization
and argumentative purposes
are not identical.

% Structural ID: VS.B01.P01.C08.S17.SS14.P003
\phantomsection
\label{para:VS.B01.P01.C08.S17.SS14.P003}

The influence of Ibn Arabi
upon Sadra should not
be used to collapse their
respective doctrines. The
Akbarian account of
self-disclosure is not
simply another expression
of substantial motion,
and Sadra's gradation
of existence is not
identical with the
differentiation of divine
names or the determinate
realities examined in
Section~\ref{sec:ch08_s16}.
The comparison is most
productive where it
identifies how different
metaphysical systems
address the relation
between the One and
the many while preserving
their distinctive explanatory
principles.\footnote{For the
relation between Ibn Arabi
and Mulla Sadra, see Sajjad
H. Rizvi, ``Philosophy and
Mysticism: Ibn 'Arabi and
Mulla Sadra,'' in
\textit{The Cambridge Companion
to Arabic Philosophy}, ed.
Peter Adamson and Richard C.
Taylor (Cambridge: Cambridge
University Press, 2005),
pp. 224--246.}

% Structural ID: VS.B01.P01.C08.S17.SS14.P004
\phantomsection
\label{para:VS.B01.P01.C08.S17.SS14.P004}

The transition from Ibn Arabi
to Sadra thus marks a
change in the comparative
question. The former
foregrounds the unity
of the Real and the
differentiated disclosure
of existence; the latter
develops a philosophical
account in which existence
itself is the fundamental
reality, differentiated
through gradation and
dynamic through substantial
motion. The comparison
must examine this change
without treating Sadra
as merely a philosophical
restatement of Ibn Arabi.

\index[name]{Ibn Arabi}
\index[concept]{Unity of existence}


% =============================================================================
% SS15 — THE STRONGEST SADRIAN CHALLENGE
% =============================================================================

\subsection{The Strongest Sadrian Challenge to Book I}
\label{subsec:ch08_s17_ss15}

% Structural ID: VS.B01.P01.C08.S17.SS15.P001
\phantomsection
\label{para:VS.B01.P01.C08.S17.SS15.P001}

The strongest Sadrian challenge
to Book I concerns the
ontological work assigned
to Identity. If existence
is the fundamental reality,
if it is the principle
of individuation, and
if its modulation explains
the differentiated actuality
of existents, what remains
for a distinct ontological
account of Identity to
explain? A Sadrian critic
may argue that Book I
risks introducing a
second principle where
the reality of existence
already performs the
relevant work.

% Structural ID: VS.B01.P01.C08.S17.SS15.P002
\phantomsection
\label{para:VS.B01.P01.C08.S17.SS15.P002}

The challenge is stronger
than the assertion that
existence and Identity
are merely different words.
It asks whether the
distinction between Being
and Identity introduces
an explanatory duplication.
If the existent is actual
through existence, and
if its individuality is
also grounded in existence,
then a separate account
of Identity may appear
to redescribe the
individuation and
determination already
explained by Sadra's
metaphysics.

% Structural ID: VS.B01.P01.C08.S17.SS15.P003
\phantomsection
\label{para:VS.B01.P01.C08.S17.SS15.P003}

A second challenge concerns
the treatment of change.
Sadra's account of
substantial motion rejects
the assumption that the
existent remains an
ontologically static
substrate beneath its
transformations. A
framework that distinguishes
Identity from the
changing determinations
and manifestations of
an existent must explain
how this distinction
avoids reintroducing
a hidden, immutable
substrate. It must
show that Identity is
not simply a name for
a substance that remains
untouched by the actual
process of existence.

% Structural ID: VS.B01.P01.C08.S17.SS15.P004
\phantomsection
\label{para:VS.B01.P01.C08.S17.SS15.P004}

A third challenge concerns
the relation between
analytical distinction
and ontological reality.
Sadra may ask whether
Book I's analytical
distinction between Being
and Identity identifies
two aspects of one
ontological actuality
or whether it risks
reifying a conceptual
distinction into an
additional ontological
principle. The distinction
must therefore demonstrate
its explanatory necessity
without asserting that
Identity exists as
a separate entity
alongside Being.

% Structural ID: VS.B01.P01.C08.S17.SS15.P005
\phantomsection
\label{para:VS.B01.P01.C08.S17.SS15.P005}

The Sadrian challenge
can be formulated as
a demand for ontological
economy: if the reality
of existence accounts
for actuality, gradation,
and individuation, Book I
must identify the precise
question that remains
unanswered. It cannot
answer by merely repeating
that Being and Identity
are distinguishable.
It must show what
the distinction enables
one to explain that
cannot be obtained
from the Sadrian account
of existence and
individuation alone.

\index[concept]{Ontological economy}


% =============================================================================
% SS16 — THE STRONGEST BOOK I CHALLENGE
% =============================================================================

\subsection{The Strongest Book I Challenge to Sadra}
\label{subsec:ch08_s17_ss16}

% Structural ID: VS.B01.P01.C08.S17.SS16.P001
\phantomsection
\label{para:VS.B01.P01.C08.S17.SS16.P001}

The strongest challenge from
Book I concerns whether
the primacy of existence
and the principle of
individuation fully explain
the existent as the
existent that it is.
Sadra's account gives
existence a fundamental
role in actuality and
individuality. Yet the
question remains whether
the existential actuality
and individuation of
an entity exhaust the
ontological meaning
of its Identity.

% Structural ID: VS.B01.P01.C08.S17.SS16.P002
\phantomsection
\label{para:VS.B01.P01.C08.S17.SS16.P002}

The distinction is not
between existence and
a hidden essence that
subsists independently
of existence. Book I
does not adopt such
an essentialist position.
Its question concerns
the relation between
Being and Identity:
Being concerns the
actuality of what exists,
while Identity concerns
the existent as the
existent that it is.
The two are analytically
distinguishable and
ontologically inseparable,
without ontological
priority of either over
the other. This relation
cannot be translated
directly into the
Sadrian distinction
between existence and
quiddity.

% Structural ID: VS.B01.P01.C08.S17.SS16.P003
\phantomsection
\label{para:VS.B01.P01.C08.S17.SS16.P003}

A further challenge concerns
the meaning of individuation.
The principle through
which an existent is
constituted as an
individual does not
necessarily explain
every question concerning
its Identity. Individuality,
numerical distinction,
and Identity are related
but non-equivalent
problems. If Sadra's
account of existence
as the principle of
individuation is to
replace the question
of Identity, it must
demonstrate that no
ontological question
remains after individuation
has been explained.

% Structural ID: VS.B01.P01.C08.S17.SS16.P004
\phantomsection
\label{para:VS.B01.P01.C08.S17.SS16.P004}

The problem becomes more
difficult in the context
of substantial motion.
If the existent is
continuously transformed
in its substantial
actuality, the account
must clarify the relation
between the changing
existential stages,
the individuality of
the subject, and the
Identity of the existent.
The fact that a process
is continuous does
not itself establish
the numerical identity
of every stage. Likewise,
the fact that substantial
change occurs does
not itself establish
that Identity has
terminated. These
questions require
an account of their
own.

% Structural ID: VS.B01.P01.C08.S17.SS16.P005
\phantomsection
\label{para:VS.B01.P01.C08.S17.SS16.P005}

Book I therefore asks
whether Sadra's metaphysics
provides a sufficient
account of ontological
specificity, rather than
whether it successfully
explains existence
and change. Its challenge
is not that Sadra has
ignored individuality
or failed to recognize
the reality of becoming.
It is that the relation
between Being, individuation,
and Identity must be
made explicit before
one can conclude that
the first two exhaust
the third.

\index[concept]{Ontological specificity}


% =============================================================================
% SS17 — THE SUBSTITUTION TEST
% =============================================================================

\subsection{The Substitution Test}
\label{subsec:ch08_s17_ss17}

% Structural ID: VS.B01.P01.C08.S17.SS17.P001
\phantomsection
\label{para:VS.B01.P01.C08.S17.SS17.P001}

The substitution test examines
whether the principal
Sadrian concepts can
replace the corresponding
concepts of Book I
without changing the
meaning or inferential
role of the claims.
The test does not assume
that concepts from
different philosophical
systems must be
incommensurable. It
asks whether a proposed
substitution preserves
the work performed
by the concept in
its original system.

% Structural ID: VS.B01.P01.C08.S17.SS17.P002
\phantomsection
\label{para:VS.B01.P01.C08.S17.SS17.P002}

The first proposed substitution
is between \textit{wujūd}
and Being. The concepts
are closely related,
but their full functions
are not automatically
identical. Sadra's
\textit{wujūd} carries
the specific commitments
of the primacy, modulation,
and dependence of
existence. Book I's
Being is examined
through its relation
to Identity and the
absence of ontological
priority between them.
The substitution is
therefore partial and
requires qualification.

% Structural ID: VS.B01.P01.C08.S17.SS17.P003
\phantomsection
\label{para:VS.B01.P01.C08.S17.SS17.P003}

The second proposed substitution
is between
\textit{tashakhkhuṣ}
and Identity. The former
concerns individuation;
the latter concerns
the existent as the
existent that it is.
Although individuation
is directly relevant
to Identity, it does
not necessarily exhaust
its ontological significance.
The substitution therefore
fails unless an argument
establishes that the
constitution of an
individual and the
Identity of that
individual are
coextensive.

% Structural ID: VS.B01.P01.C08.S17.SS17.P004
\phantomsection
\label{para:VS.B01.P01.C08.S17.SS17.P004}

The third proposed substitution
is between
\textit{tashkīk al-wujūd}
and Invariance. The
gradation of existence
concerns differentiated
degrees of existential
actuality. Invariance
concerns what remains
invariant under a
specified transformation,
within a specified
domain and respect.
The two concepts
address different
questions. A graded
existence may undergo
substantial change,
and an invariant
may be specified
without being a
degree of existence.
The substitution
therefore fails.

% Structural ID: VS.B01.P01.C08.S17.SS17.P005
\phantomsection
\label{para:VS.B01.P01.C08.S17.SS17.P005}

The fourth proposed substitution
is between substantial
motion and Expression.
Substantial motion
concerns the dynamic
actuality of the
substance itself.
Expression concerns
a distinct ontological
relation developed
in Book I. Both may
be relevant to the
dynamism of an existent,
but they do not
have the same scope
or explanatory role.
The substitution
fails without a
further argument
that establishes
their equivalence.

% Structural ID: VS.B01.P01.C08.S17.SS17.P006
\phantomsection
\label{para:VS.B01.P01.C08.S17.SS17.P006}

The test therefore yields
a differentiated result.
Sadra's doctrines provide
powerful accounts of
existence, gradation,
individuation, and
substantial change.
They illuminate several
problems central to
Book I. Yet no
one-to-one substitution
has been established
between those doctrines
and Being, Identity,
Invariance, Expression,
or Continuity. The
relation is one of
philosophical proximity
and critical engagement,
not conceptual identity.

\index[concept]{Substitution test}


% =============================================================================
% SS18 — COMPARATIVE IMPLICATIONS
% =============================================================================

\subsection{Comparative Implications for Being and Identity}
\label{subsec:ch08_s17_ss18}

% Structural ID: VS.B01.P01.C08.S17.SS18.P001
\phantomsection
\label{para:VS.B01.P01.C08.S17.SS18.P001}

The comparison establishes
that Sadra's metaphysics
offers a substantial
account of the actuality
of existents. By treating
existence as fundamental,
he provides an ontological
basis for the reality
of contingent beings.
Through gradation, he
explains how one existential
reality can be differentiated
without being reduced
to a plurality of
independent existential
principles. Through
individuation, he
connects existence
to the constitution
of the individual.
Through substantial
motion, he places
becoming within the
actuality of existence.

% Structural ID: VS.B01.P01.C08.S17.SS18.P002
\phantomsection
\label{para:VS.B01.P01.C08.S17.SS18.P002}

These achievements provide
three important challenges
for Book I. First,
the account of Being
must preserve the
reality of existence
and avoid reducing
it to a conceptual
or descriptive category.
Second, the account
of Identity must
explain its relation
to individuation
without identifying
the two. Third,
the account of
change must address
the possibility
that substantial
transformation occurs
within the actuality
of an existent,
rather than merely
in its accidental
determinations or
manifestations.

% Structural ID: VS.B01.P01.C08.S17.SS18.P003
\phantomsection
\label{para:VS.B01.P01.C08.S17.SS18.P003}

The comparison also
clarifies the limits
of the apparent opposition
between Sadra and
Book I. Sadra's
primacy of existence
is directed toward
the relation between
existence and quiddity.
Book I's absence
of ontological priority
concerns the relation
between Being and
Identity. Since quiddity
and Identity are
not equivalent, the
two positions cannot
be treated as direct
contradictions. A
genuine conflict
would have to be
established through
a further argument
showing that Sadra's
account of existence
excludes the analytical
distinction between
Being and Identity.

% Structural ID: VS.B01.P01.C08.S17.SS18.P004
\phantomsection
\label{para:VS.B01.P01.C08.S17.SS18.P004}

The following table
summarizes the principal
Sadrian doctrines and
the limits of their
comparison with the
foundational concepts
of Book I. It is
a comparative guide,
not a reduction
of either system
to the terminology
of the other.

% =============================================================================
% Project : THE_PHILOSOPHY_OF_VIABILITY
% Book    : Book I — Ontology of Viability
% Book ID : VS.B01
% Part    : Part I — Identity and Manifestation
% Part ID : VS.B01.P01
% Chapter : Chapter 08 — Comparative Philosophical Examination
% Chapter ID : VS.B01.P01.C08
%
% File    :
%   4_ASSETS/tables/
%   VS.B01.P01.C08.S17.SS18.P004.TL01.sadrian_doctrines_and_comparative_limits.tex
%
% Asset   : TL01 — Sadrian Doctrines and Comparative Limits
% Structural Identifier :
%   VS.B01.P01.C08.S17.SS18.P004.TL01
% Asset Type : TL
%
% Purpose :
%   Summarize the principal Sadrian doctrines, their logical
%   form and epistemic status in the present reconstruction,
%   and their comparative limits in relation to Book I.
%
% =============================================================================

\begin{table}[htbp]
\centering
\footnotesize
\renewcommand{\arraystretch}{1.3}
\caption{Sadrian doctrines, claim status, and comparative limits}
\label{tab:ch08_s17_sadrian_doctrines}

\begin{tabular}{
    p{0.13\textwidth}
    p{0.10\textwidth}
    p{0.16\textwidth}
    p{0.17\textwidth}
    p{0.17\textwidth}
    p{0.17\textwidth}
}
\toprule
\textbf{Doctrine or claim}
&
\textbf{Logical form}
&
\textbf{Epistemic status}
&
\textbf{Sadrian function}
&
\textbf{Book I point of contact}
&
\textbf{Comparative limit}
\\
\midrule

\textit{Aṣālat al-wujūd}
&
Thesis
&
Conditionally Supported Claims
&
Existence is the fundamental ontological reality.
&
Being and the actuality of the existent.
&
Primacy over quiddity is not priority over Identity.
\\

\textit{Tashkīk al-wujūd}
&
Thesis
&
Conditionally Supported Claims
&
Existence is realized through differentiated degrees.
&
Unity, multiplicity, and differentiation.
&
Gradation is not Invariance or the realm--hub architecture.
\\

\textit{Tashakhkhuṣ}
&
Thesis
&
Conditionally Supported Claims
&
Existence functions as the principle of individuation.
&
The ontological specificity of an existent.
&
Individuation does not automatically exhaust Identity.
\\

\textit{Al-ḥaraka al-jawhariyya}
&
Thesis
&
Conditionally Supported Claims
&
Substantial actuality is itself dynamic.
&
Change and continuity of the existent.
&
Substantial motion is not Expression or Manifestation.
\\

Knowledge by presence and the union of intellect and intelligible
&
Thesis
&
Conditionally Supported Claims
&
Knowing is examined through presence and intellectual union.
&
The relation between actuality and knowing.
&
These doctrines do not make Being or Identity dependent upon a knower.
\\

Non-substitutability of Sadrian and Book I concepts
&
Thesis
&
Argued/Derived Claims
&
The concepts perform related but non-identical explanatory work.
&
Being, Identity, Invariance, and Expression.
&
The comparison establishes proximity, not equivalence.
\\

The relation between substantial becoming and Identity
&
Open Problem
&
Unresolved/Open Problems
&
The continuity and individuality of the changing existent require further interpretation.
&
Identity through substantial transformation.
&
No final resolution is adopted in this section.
\\

\bottomrule
\end{tabular}
\end{table}

% =============================================================================
% END OF TABLE TL01
% =============================================================================

% Structural ID: VS.B01.P01.C08.S17.SS18.P005
\phantomsection
\label{para:VS.B01.P01.C08.S17.SS18.P005}

The principal implication
is that the Sadrian
account should be
treated as a serious
alternative explanation
of existential actuality,
individuation, and
change. It does
not establish that
Identity is reducible
to existence, nor
does it establish
that Book I's
distinction is
unnecessary. The
comparison instead
places a precise
burden upon Book I:
it must demonstrate
the explanatory
significance of
Identity without
turning it into
a second substance,
a quiddity independent
of existence, or
a synonym for
individuation.

\index[concept]{Being and Identity}
\index[concept]{Identity and individuation}


% =============================================================================
% SS19 — OPEN PROBLEMS AND CLAIM STATUS
% =============================================================================

\subsection{Open Problems and Claim Status}
\label{subsec:ch08_s17_ss19}

% Structural ID: VS.B01.P01.C08.S17.SS19.P001
\phantomsection
\label{para:VS.B01.P01.C08.S17.SS19.P001}

Several questions remain
unresolved. First,
does Sadra's principle
of individuation
through existence
provide a sufficient
account of the
ontological specificity
of an existent,
or does it leave
a distinct question
concerning Identity?
Second, how should
the relation between
the continuity of
substantial becoming
and the individuality
of the existent
be formulated? Third,
does the account
of the soul provide
a model of continuity
that can be extended
beyond human personhood,
or is its explanatory
scope restricted
to the specific
metaphysical status
of the soul?

% Structural ID: VS.B01.P01.C08.S17.SS19.P002
\phantomsection
\label{para:VS.B01.P01.C08.S17.SS19.P002}

Further questions concern
the interpretation
of gradation, the
relation between
unity and multiplicity,
the ontological
significance of
successive existential
stages, and the
relation between
existence as a
principle of individuation
and Identity as
examined in Book I.
These questions
cannot be settled
merely by identifying
similar terms or
by assuming that
one system contains
the complete answer
to the problems
of the other.

% Structural ID: VS.B01.P01.C08.S17.SS19.P003
\phantomsection
\label{para:VS.B01.P01.C08.S17.SS19.P003}

The claim-status record
must distinguish the
historical reconstruction
of Sadra's doctrines
from the comparative
theses advanced in
this section. The
primacy of existence,
gradation, individuation,
and substantial motion
are reconstructed
as doctrines within
Sadra's system. Their
historical attribution
is supported by
primary-text references
and secondary scholarship.
The claim that these
doctrines are not
substitutable without
qualification for
the concepts of
Book I is a
comparative thesis,
not a foundational
commitment of Sadra
or of Book I.

% Structural ID: VS.B01.P01.C08.S17.SS19.P004
\phantomsection
\label{para:VS.B01.P01.C08.S17.SS19.P004}

The distinction between
the functional or
logical form of
a claim and
its epistemic or
justificatory status
must remain explicit.
A thesis may be
conditionally supported;
an argument may
remain incomplete;
a historical attribution
may be well-supported
without being adopted
as a foundational
commitment. The
comparison therefore
does not turn
Sadra's metaphysical
principles into
principles of Book I.
Nor does it
treat an unresolved
question as though
it had already
been answered.

% Structural ID: VS.B01.P01.C08.S17.SS19.P005
\phantomsection
\label{para:VS.B01.P01.C08.S17.SS19.P005}

The section's principal
comparative thesis
is that Sadra's
metaphysics offers
a powerful account
of existential actuality,
gradation, individuation,
and substantial change,
while the relation
between these doctrines
and the ontological
meaning of Identity
remains open to
further argument.
The thesis is
conditionally supported
by the reconstruction
and comparison
developed here.
It does not
establish the
superiority of
either system,
nor does it
settle the
ultimate ground
of Identity.

% Structural ID: VS.B01.P01.C08.S17.SS19.P006
\phantomsection
\label{para:VS.B01.P01.C08.S17.SS19.P006}

The open problems
are not defects
to be concealed
by a premature
verdict. They
identify the precise
points at which
the comparison
requires further
textual and
philosophical work.
The relation
between existential
actuality, individuation,
substantial becoming,
and Identity
remains a central
question for
the ontology
developed in
Part I.

\index[concept]{Claim status}
\index[concept]{Open ontological problems}


% =============================================================================
% SS20 — TRANSITION TO NAGARJUNA
% =============================================================================

\subsection{Transition to Nagarjuna}
\label{subsec:ch08_s17_ss20}

% Structural ID: VS.B01.P01.C08.S17.SS20.P001
\phantomsection
\label{para:VS.B01.P01.C08.S17.SS20.P001}

The examination of Sadra
has centered on
the primacy and
gradation of existence,
the principle of
individuation, and
substantial motion.
His metaphysics offers
an account in
which existential
actuality is
fundamental, multiplicity
is articulated
through degrees
of existence,
and becoming
is located
within the
substance of
the existent itself.
The comparison
has shown that
these doctrines
provide a serious
challenge to
any account
that distinguishes
Being and Identity
without explaining
the ontological
work performed
by that distinction.

% Structural ID: VS.B01.P01.C08.S17.SS20.P002
\phantomsection
\label{para:VS.B01.P01.C08.S17.SS20.P002}

The next comparison
moves to a
different philosophical
orientation. Nagarjuna's
analysis of
dependent origination
and emptiness
questions the
assumption that
entities possess
an independently
established intrinsic
nature. The
transition therefore
moves from
a metaphysics
that gives
existence a
fundamental and
modulated reality
to an examination
of whether
the concepts
of intrinsic
nature, dependence,
and identity
can be sustained
without reifying
the existent.

% Structural ID: VS.B01.P01.C08.S17.SS20.P003
\phantomsection
\label{para:VS.B01.P01.C08.S17.SS20.P003}

This transition
does not assume
that Nagarjuna's
critique of
intrinsic nature
directly refutes
Sadra's metaphysics,
nor that
emptiness is
a synonym
for non-existence.
It introduces
a distinct
problem concerning
the conditions
under which
an existent
can be
described as
dependent, determinate,
and identifiable.
The next
section must
reconstruct Nagarjuna's
arguments on
their own terms
before examining
their implications
for the ontological
distinctions of
Book I.

% =============================================================================
% END OF SECTION 17 OF CHAPTER 08
% =============================================================================